\SourceLeaf{263}%
nicht wahrnehmbar gelten. Im größten der Dreiecke, die wir in den vergangenen Jahren vermessen haben, nämlich zwischen den Punkten Hohenhagen, Brocken und Inselsberg, hatte die Winkelsumme beispielsweise einen Überschuss \(=14''85348\). Die Rechnung ergab die folgenden Korrekturen der Winkel.
\begin{center}
\begin{tabular}{ll}
Hohenhagen \(.\ .\ .\ .\ .\)&\(-4''95113\)\\
Brocken \(.\ .\ .\ .\ .\ .\)&\(-4{,}95104\)\\
Inselsberg\(.\ .\ .\ .\ .\ .\)&\(-4{,}95131\)
\end{tabular}
\end{center}
\Article{29}
\indent Zum Schluss wollen wir noch den Flächeninhalt eines Dreiecks auf einer gekrümmten Fläche mit dem Flächeninhalt eines geradlinigen Dreiecks vergleichen, dessen Seiten \(a,b,c\) sind. Den Flächeninhalt des letzteren bezeichnen wir mit \(\sigma^*\) \(=\tfrac{1}{2}bc\sin A^*=\tfrac{1}{2}ac\sin B^*=\tfrac{1}{2}ab\sin C^*\)

Bis zu Größen vierter Ordnung gilt
\[
\sin A^*=\sin A-\tfrac{1}{12}\sigma\cos A\mathord{.}(2\alpha+\beta+\gamma)
\]
oder mit derselben Genauigkeit
\[
\sin A=\sin A^*\mathord{.}(1+\tfrac{1}{24}bc\cos A\mathord{.}(2\alpha+\beta+\gamma))
\]
Setzen wir diesen Wert in Formel [9] ein, so erhalten wir bis zu Größen sechster Ordnung
\[
\begin{aligned}
\sigma=\tfrac{1}{2}bc\sin A^*\mathord{.}\{1&+\tfrac{1}{120}\alpha(3bb+3cc-2bc\cos A)\\
&+\tfrac{1}{120}\beta(3bb+4cc-4bc\cos A)\\
&+\tfrac{1}{120}\gamma(4bb+3cc-4bc\cos A)\}
\end{aligned}
\]
oder mit derselben Genauigkeit
\[
\begin{aligned}
\sigma=\sigma^*\{1&+\tfrac{1}{120}\alpha(aa+2bb+2cc)+\tfrac{1}{120}\beta(2aa+bb+2cc)\\
&+\tfrac{1}{120}\gamma(2aa+2bb+cc)\}
\end{aligned}
\]
Für eine Kugelfläche nimmt diese Formel folgende Gestalt an
\[
\sigma=\sigma^*(1+\tfrac{1}{24}\alpha(aa+bb+cc))
\]
Wie sich leicht bestätigen lässt, kann man mit derselben Genauigkeit auch die folgende Formel dafür verwenden
\[
\sigma=\sigma^*\BareRoot{\frac{\sin A\mathord{.}\sin B\mathord{.}\sin C}{\sin A^*\mathord{.}\sin B^*\mathord{.}\sin C^*}}
\]
Wendet man dieselbe Formel auf Dreiecke auf einer nicht kugelförmigen gekrümmten Fläche an, so ist der Fehler im Allgemeinen von fünfter Ordnung. Er ist jedoch bei allen Dreiecken, wie sie sich auf der Erdoberfläche vermessen lassen, nicht wahrnehmbar.
\par\medskip\centerline{\rule{.75in}{.3pt}}
